% భాగం: సమితులు, సంబంధాలు, ప్రమేయాలు

\documentclass[../../include/open-logic-part]{subfiles}

\begin{document}

\olpart{sfr}{సమితులు, సంబంధాలు, ప్రమేయాలు}

\begin{editorial}
  సరళ సమితి సిద్ధాంత భాగంలోని మొదటి నాలుగు అధ్యాయాలను ఈ దస్త్రం వాటి తొలి రూపంలో చేర్చుతుంది. ఆ రూపంలో భావాలను నెమ్మదిగా పరిచయం చేస్తారు. పూర్తి భాగాన్ని \verb|sets-functions-relations-complete.tex| చేర్చుతుంది. ఇక్కడి పాఠ్యంతో సరళ సమితి సిద్ధాంతపు మౌలిక విషయాలకు తగినంత సమగ్ర పరిచయం లభిస్తుంది. ఈ విషయాలు విద్యార్థులకు ముందే తెలుసని అనుకోలేకపోతే, ఈ పాఠ్యాన్ని, కనీసం సమితులు, ప్రమేయాలు, సంబంధాలకు సంబంధించిన భాగాన్నైనా మళ్లీ పరిశీలిస్తూ పాఠ్యక్రమాన్ని మొదలుపెట్టడం మంచిది. అలా చేస్తే ఇతర తార్కిక విభాగాలకు అవసరమైన గణిత సంకేతాలను వాడటంలో విద్యార్థులందరికీ ప్రాథమిక నైపుణ్యం వస్తుంది.
\end{editorial}

\olimport[sets]{sets}

\olimport[relations]{relations}

\olimport[functions]{functions}

\olimport[size-of-sets]{size-of-sets}

\OLEndPartHook
\end{document}
